\documentclass[UTF8,10pt,oneside]{ctexbook}

\usepackage[a4paper,margin=2.0cm]{geometry}
\usepackage{fontspec}
\usepackage{xcolor}
\usepackage{hyperref}
\usepackage{xurl}
\usepackage{graphicx}
\usepackage{amsmath,amssymb}
\usepackage{booktabs}
\usepackage{longtable}
\usepackage{tabularx}
\usepackage{array}
\usepackage{multirow}
\usepackage{enumitem}
\usepackage{fancyhdr}
\usepackage{setspace}
\usepackage{float}
\usepackage{caption}
\usepackage[most]{tcolorbox}
\usepackage{tikz}
\usepackage{pgfplots}
\usetikzlibrary{arrows.meta,positioning,shapes.geometric,shapes.misc,fit,calc}
\pgfplotsset{compat=1.18}

\IfFontExistsTF{PingFang SC}{
  \setCJKmainfont{PingFang SC}
}{
  \IfFontExistsTF{Hiragino Sans GB}{
    \setCJKmainfont{Hiragino Sans GB}
  }{}
}
\IfFontExistsTF{Menlo}{
  \setmonofont{Menlo}
}{}

\definecolor{DeepBlue}{HTML}{1F4E79}
\definecolor{SoftBlue}{HTML}{EEF5FB}
\definecolor{SoftGreen}{HTML}{EDF7F0}
\definecolor{SoftOrange}{HTML}{FFF3E6}
\definecolor{SoftGray}{HTML}{F5F7FA}
\definecolor{DarkText}{HTML}{1D2733}
\definecolor{AccentRed}{HTML}{C0392B}
\definecolor{AccentGreen}{HTML}{2E7D32}

\hypersetup{
  colorlinks=true,
  linkcolor=DeepBlue,
  urlcolor=DeepBlue,
  pdftitle={从 PDF 到分析结论：第一性原理教学讲义},
  pdfauthor={Codex},
  pdfsubject={Graphiti 本地案例教学资料}
}

\setstretch{1.05}
\setlength{\parindent}{2em}
\setlength{\parskip}{0.12em}
\setlist[itemize]{leftmargin=2.2em}
\setlist[enumerate]{leftmargin=2.5em}

\pagestyle{fancy}
\fancyhf{}
\fancyhead[L]{从 PDF 到分析结论：第一性原理讲义}
\fancyhead[R]{武汉大学案例}
\fancyfoot[C]{\thepage}
\setlength{\headheight}{14pt}

\tcbset{
  enhanced,
  breakable,
  boxrule=0.5pt,
  arc=2mm
}

\newtcolorbox{definitionbox}[1][]{
  colback=SoftBlue,
  colframe=DeepBlue,
  title=#1,
  fonttitle=\bfseries
}

\newtcolorbox{casebox}[1][]{
  colback=SoftOrange,
  colframe=orange!60!black,
  title=#1,
  fonttitle=\bfseries
}

\newtcolorbox{takeawaybox}[1][]{
  colback=SoftGreen,
  colframe=AccentGreen,
  title=#1,
  fonttitle=\bfseries
}

\newtcolorbox{warningbox}[1][]{
  colback=white,
  colframe=AccentRed,
  title=#1,
  fonttitle=\bfseries
}

\newcommand{\stage}[1]{\textbf{#1}}
\newcommand{\code}[1]{\nolinkurl{#1}}

\captionsetup{font=small,labelfont=bf}

\begin{document}

\frontmatter

\begin{titlepage}
\centering
\vspace*{2.5cm}
{\Huge\bfseries 从 PDF 到分析结论\\[0.4em]第一性原理教学讲义\par}
\vspace{1cm}
{\Large 基于 Graphiti 本地部署与武汉大学舆情案例的系统推导\par}
\vspace{1.5cm}

\begin{tcolorbox}[width=0.88\textwidth,colback=SoftBlue,colframe=DeepBlue]
\centering
\large
这不是一份“功能说明书”，而是一份解释系统为何这样设计、\\
每一层中间表示在解决什么问题、以及为什么最终能从 PDF 走到分析结论的教学讲义。
\end{tcolorbox}

\vfill

\begin{tabular}{rl}
案例输入文档： & 武汉大学品牌声誉深度分析报告.pdf \\
模拟问题： & 如果武汉大学发布了对肖某某处分的撤销公告，大众舆情走向会怎么样 \\
系统形态： & Flask Backend + Graphiti Sidecar + Kuzu + OASIS 双平台模拟 \\
运行模型： & gpt-5.4-mini \\
Embedding： & BAAI/bge-m3 \\
编译方式： & XeLaTeX \\
日期： & 2026-03-29 \\
\end{tabular}

\vfill
{\large Codex 依据当前代码、配置与真实日志编写\par}
\end{titlepage}

\chapter*{摘要}

这份讲义试图回答一个看似简单、实则非常容易被误解的问题：为什么一份 PDF 文档，经过一套本地系统处理之后，能够变成一份带有趋势判断、角色分化、舆情节奏和风险提示的分析报告？

如果只从表面看，这个系统像是在做“文档理解 + 大模型生成”。但从第一性原理看，它做的事情更接近于：先把原始文档压缩成可检索的世界模型，再把世界模型变成可发声的主体集合，再让主体在一个有时间和平台约束的仿真环境中演化，最后让报告智能体回到这个世界模型和仿真世界里进行二阶推理。

因此，这套系统的本质并不是“把 PDF 喂给模型，让模型直接写答案”，而是分三层连续构造：

\begin{enumerate}
  \item 构造\textbf{可计算文本}；
  \item 构造\textbf{可检索、可复用、可被推演的中间世界模型}；
  \item 构造\textbf{基于世界模型的多主体演化与分析生成}。
\end{enumerate}

本讲义以武汉大学案例作为真实教学样本，直接使用本次运行中的正式参数和正式结果：

\begin{itemize}
  \item 原始 PDF 文件大小约 2.89MB，抽取后文本长度为 20862 字符；
  \item chunk 参数为 6500 / 650，只切成 4 个 chunk；
  \item 图谱规模为 64 个节点、80 条边；
  \item 从图谱中筛出并生成了 59 个 Agent 配置；
  \item 双平台仿真共执行 216 个动作；
  \item 报告阶段本地搜索调用 21 次，其中 18 次命中为 0；
  \item 全链路从项目创建到最终报告完成约 939.3 秒，即 15 分 39 秒。
\end{itemize}

这些数据非常适合作为教学材料，因为它们同时展示了系统为什么能工作，也展示了它为什么会慢、为什么会失真，以及为什么参数调优永远是在速度、召回和表达力之间做权衡。

\tableofcontents

\mainmatter

\chapter{问题定义：这到底是不是一个摘要问题}

\section{表面任务与真实任务}

很多人第一次看到这个案例时，会把问题理解成：“给定一份 PDF，再给定一个问题，让大模型输出一份分析报告。”这当然不是完全错，但它省略了最关键的部分：\textbf{报告中的许多内容并不直接存在于原文里，而是来自对原文结构的重建、对主体关系的外推、对多角色反应的模拟，以及对这些模拟结果的归纳。}

换句话说，系统面对的不是一个简单的“抽取事实”任务，而是一个混合任务：

\begin{itemize}
  \item 一部分是\textbf{事实提取}：原文里提到了哪些主体、关系、事件、平台和观点；
  \item 一部分是\textbf{结构重建}：这些事实如何被组织成一个可检索的网络；
  \item 一部分是\textbf{行为建模}：如果这些主体都能发声，他们会如何响应一个新的事件假设；
  \item 一部分是\textbf{分析写作}：如何把上面三部分重新组织成对人可读的结论。
\end{itemize}

\begin{definitionbox}[关键定义：什么叫“从 PDF 到分析”]
所谓“从 PDF 到分析”，不是把一个文件变成一段更短的文字，而是把一个非结构化信息载体，逐层转译为可被检索、可被推演、可被审计的中间表示，最后再转译为面向人的解释性输出。
\end{definitionbox}

\section{为什么不能直接把 PDF 扔给大模型}

从第一性原理看，直接把 PDF 扔给大模型有四个硬伤。

\paragraph{第一，输入太原始。}
PDF 是阅读格式，不是知识格式。它混合了版式、页面、段落、脚注和跨页断裂，对人友好，但对后续计算并不友好。

\paragraph{第二，信息不可复用。}
如果每次都直接把整份 PDF 和问题丢给模型，那么每一次回答都是一次重新消费全部上下文，成本高，而且中间结果不可复用。

\paragraph{第三，信息不可检索。}
系统如果想回答“哪些角色会支持”“哪个平台先引爆”“谁会把话题带向高校治理”，就需要一个能按实体、关系、边事实来索引的信息结构。原始 PDF 做不到这一点。

\paragraph{第四，信息不可推演。}
“如果武汉大学发布撤销公告”是一个反事实问题。原始 PDF 提供的是历史描述，而不是主体层面的演化机制。要回答反事实问题，必须先构造可发声的主体集合和它们之间的影响网络。

\begin{takeawaybox}[本章核心结论]
这套系统解决的不是“摘要”问题，而是“如何把非结构化文档转换成一个可被分析、可被搜索、可被模拟的中间世界模型”问题。报告只是这个中间世界模型的一个最终输出视图。
\end{takeawaybox}

\chapter{第一性原理：为什么必须有中间世界模型}

\section{信息处理的三层目标}

如果把这类系统抽象成信息加工管线，那么它至少要满足三个目标：

\begin{enumerate}
  \item \textbf{压缩}：把冗长文本压成较短但仍保留结构的信息表示；
  \item \textbf{索引}：让系统能够按主题、实体、关系和问题进行局部访问，而不是每次全文重读；
  \item \textbf{推演}：让系统能够不只回答“原文说了什么”，还回答“在这个世界里，若新增一个事件，会发生什么”。
\end{enumerate}

这三个目标共同决定了为什么中间表示不能省略。

\section{从载体到世界模型}

我们可以把系统的中间表示理解成一层层“再表达”：

\[
D \xrightarrow{P} T \xrightarrow{C} \{c_i\} \xrightarrow{O} \mathcal{O}
\xrightarrow{G} \mathcal{G} \xrightarrow{A} \mathcal{A}
\xrightarrow{\text{sim}} \mathcal{S} \xrightarrow{R} \text{Report}
\]

其中：

\begin{itemize}
  \item \(D\) 表示原始 PDF；
  \item \(T\) 表示被抽取并清洗后的纯文本；
  \item \(\{c_i\}\) 表示文本块集合；
  \item \(\mathcal{O}\) 表示本体，即“系统允许哪些主体和关系存在”；
  \item \(\mathcal{G}\) 表示图谱，即中间世界模型；
  \item \(\mathcal{A}\) 表示 Agent 集合；
  \item \(\mathcal{S}\) 表示社会仿真轨迹；
  \item Report 是面向人的最终叙事。
\end{itemize}

\begin{definitionbox}[中间世界模型]
在本讲义中，“中间世界模型”特指图谱 \(\mathcal{G}\) 及其可派生的主体集合 \(\mathcal{A}\)。它不是全文复制，也不是最终结论，而是介于两者之间、能够支持检索和推演的结构层。
\end{definitionbox}

\section{每一层都在保留什么、丢弃什么、获得什么}

\begin{longtable}{p{2.2cm}p{3.6cm}p{3.6cm}p{4.2cm}}
\toprule
层级 & 保留的东西 & 丢弃的东西 & 新获得的能力 \\
\midrule
PDF $\rightarrow$ 文本 & 原始语义内容 & 页面布局、视觉样式 & 可被程序读取 \\
文本 $\rightarrow$ chunk & 局部语义片段 & 长距离连续性的一部分 & 可分批处理、可局部写入 \\
chunk $\rightarrow$ 本体 & 角色类别、关系类别 & 无关细枝末节 & 约束系统“该提取什么” \\
本体 + chunk $\rightarrow$ 图谱 & 实体、关系、事实 & 原文冗余表达 & 可检索、可索引、可统计 \\
图谱 $\rightarrow$ Agent & 可发声主体及其属性 & 与主体无关的孤立结构 & 可做行为模拟 \\
Agent + 时间 $\rightarrow$ 仿真轨迹 & 平台行为、观点演化 & 静态视角 & 可回答反事实问题 \\
图谱 + 仿真 $\rightarrow$ 报告 & 事实、视角、趋势判断 & 底层过程细节 & 面向人的解释与决策支持 \\
\bottomrule
\end{longtable}

\begin{figure}[H]
\centering
\begin{tikzpicture}[
  node distance=7mm,
  level/.style={rounded rectangle, draw=DeepBlue, thick, fill=SoftBlue, minimum width=10.5cm, minimum height=1cm, align=center},
  gain/.style={rounded rectangle, draw=AccentGreen, fill=SoftGreen, minimum width=5.4cm, minimum height=0.9cm, align=center},
  arrow/.style={-Latex, thick, DeepBlue}
]
\node[level] (pdf) {PDF：可读但不可直接计算的原始载体};
\node[level, below=of pdf] (text) {文本：去版式后的语义序列};
\node[level, below=of text] (chunk) {Chunk：适合写入和召回的局部片段};
\node[level, below=of chunk] (graph) {图谱：可索引、可复用、可统计的世界模型};
\node[level, below=of graph] (agent) {Agent：可在平台上发声和互动的主体};
\node[level, below=of agent] (report) {报告：对世界模型和仿真结果的二阶解释};

\node[gain, right=12mm of pdf] (g1) {获得：可解析};
\node[gain, right=12mm of text] (g2) {获得：可分块};
\node[gain, right=12mm of chunk] (g3) {获得：可结构化};
\node[gain, right=12mm of graph] (g4) {获得：可推演};
\node[gain, right=12mm of agent] (g5) {获得：可采访};
\node[gain, right=12mm of report] (g6) {获得：可解释};

\draw[arrow] (pdf) -- (text);
\draw[arrow] (text) -- (chunk);
\draw[arrow] (chunk) -- (graph);
\draw[arrow] (graph) -- (agent);
\draw[arrow] (agent) -- (report);
\end{tikzpicture}
\caption{信息在系统中不是简单流动，而是在每一层发生“压缩 + 重表达”}
\end{figure}

\begin{takeawaybox}[本章核心结论]
如果没有中间世界模型，系统就只能反复重读全文，无法稳定检索、无法做多主体仿真，也无法把报告生成建立在可审计的结构层之上。
\end{takeawaybox}

\chapter{系统全链路总览}

\section{真实案例的系统结构}

这次案例运行使用的是一套本地化但并非全离线的结构：

\begin{itemize}
  \item 后端：Flask；
  \item 图谱存储与抽取：本地 Graphiti Sidecar + Kuzu；
  \item Embedding：本地 \code{BAAI/bge-m3}；
  \item 生成式大模型：远端 OpenAI-compatible 接口，模型为 \code{gpt-5.4-mini}；
  \item 仿真环境：OASIS 双平台并行模拟，平台为 Twitter 和 Reddit；
  \item 报告生成：ReportAgent，采用 ReACT 工具调用范式。
\end{itemize}

\begin{casebox}[本次案例正式配置]
\begin{tabularx}{\textwidth}{p{4.2cm}X}
\toprule
配置项 & 实际值 \\
\midrule
输入 PDF & 武汉大学品牌声誉深度分析报告.pdf \\
模拟问题 & 如果武汉大学发布了对肖某某处分的撤销公告，大众舆情走向会怎么样 \\
图谱后端 & graphiti \\
Graphiti 服务地址 & \code{http://127.0.0.1:8011} \\
运行模型 & \code{gpt-5.4-mini} \\
Embedding 模型 & \code{BAAI/bge-m3} \\
默认 chunk size & 6500 \\
默认 overlap & 650 \\
Graph build batch size & 2 \\
\bottomrule
\end{tabularx}
\end{casebox}

\section{从代码模块看职责分工}

\begin{table}[H]
\centering
\small
\begin{tabularx}{\textwidth}{p{4.3cm}X}
\toprule
模块 & 主要职责 \\
\midrule
\code{backend/app/utils/file\_parser.py} & 从 PDF / MD / TXT 提取文本，并做 chunk 切分 \\
\code{backend/app/services/text\_processor.py} & 统一文本预处理、分块和统计 \\
\code{backend/app/services/ontology\_generator.py} & 基于模拟问题与文档内容生成实体/关系本体 \\
\code{backend/app/services/graph\_builder.py} & 创建图谱、设置本体、分批写入 episode、等待构图完成 \\
\code{graphiti\_service/service.py} & 本地图谱侧车服务，负责 LLM 抽取、embedding、Kuzu 存储 \\
\code{backend/app/services/zep\_entity\_reader.py} & 从图谱读取节点和边，过滤出可用主体 \\
\code{backend/app/services/simulation\_config\_generator.py} & 生成时间配置、事件配置、Agent 配置和平台配置 \\
\code{backend/scripts/run\_parallel\_simulation.py} & 运行双平台并行仿真，并保持 interview 命令模式 \\
\code{backend/app/services/report\_agent.py} & 按章节执行 ReACT 式报告生成 \\
\code{backend/app/services/zep\_tools.py} & 搜索图谱、采访 Agent、提供报告工具接口 \\
\bottomrule
\end{tabularx}
\caption{从代码角度看本系统的主路径}
\end{table}

\begin{figure}[H]
\centering
\begin{tikzpicture}[
  box/.style={rounded rectangle, draw=DeepBlue, thick, fill=SoftBlue, minimum width=3.1cm, minimum height=1cm, align=center},
  box2/.style={rounded rectangle, draw=orange!70!black, thick, fill=SoftOrange, minimum width=3.1cm, minimum height=1cm, align=center},
  arrow/.style={-Latex, thick, DeepBlue}
]
\node[box] (pdf) {PDF};
\node[box, right=8mm of pdf] (text) {文本抽取};
\node[box, right=8mm of text] (chunk) {Chunk 切分};
\node[box, right=8mm of chunk] (onto) {本体生成};
\node[box, below=10mm of text] (graph) {图谱构建};
\node[box, left=8mm of graph] (entity) {实体筛选};
\node[box, right=8mm of graph] (agent) {Agent 配置};
\node[box2, below=11mm of graph] (sim) {双平台仿真};
\node[box2, right=8mm of sim] (report) {报告生成};

\draw[arrow] (pdf) -- (text);
\draw[arrow] (text) -- (chunk);
\draw[arrow] (chunk) -- (onto);
\draw[arrow] (onto.south) -- ++(0,-5mm) -| (graph.north);
\draw[arrow] (graph) -- (entity);
\draw[arrow] (graph) -- (agent);
\draw[arrow] (entity) |- (sim);
\draw[arrow] (agent) -- (sim);
\draw[arrow] (graph) |- (report);
\draw[arrow] (sim) -- (report);
\end{tikzpicture}
\caption{本案例的真实系统主路径}
\end{figure}

\begin{takeawaybox}[本章核心结论]
系统不是“一次调用一个大模型”，而是多个子系统串联：文档解析、结构化抽取、图谱持久化、主体配置、行为仿真、报告生成。每一层都在给下一层提供更高密度、更可操作的中间表示。
\end{takeawaybox}

\chapter{第一步：PDF 如何变成可计算文本}

\section{PDF 只是阅读容器，不是知识对象}

从工程角度看，PDF 最难的地方不在“读不到”，而在“读到了也未必是你想要的语义顺序”。一个面向阅读的页面结构，在程序里经常表现为：

\begin{itemize}
  \item 跨页句子被切开；
  \item 标题、页眉、页脚与正文混在一起；
  \item 原本属于同一个逻辑段落的内容在提取后变成多行；
  \item 表格、图片说明、引用框可能失去布局关系。
\end{itemize}

因此，第一步绝不是“拿到字符串”那么简单，而是把阅读载体转换成后续处理能够接受的语义序列。

\section{本项目如何抽取文本}

本项目的 \code{FileParser} 对 PDF 使用 \code{PyMuPDF} 逐页读取文本，然后把每页有内容的文本拼接起来。抽取完成后，\code{TextProcessor.preprocess\_text} 会做三件很基础但必要的事情：

\begin{enumerate}
  \item 统一换行符；
  \item 压缩连续空行；
  \item 去掉每一行首尾空白。
\end{enumerate}

这一步看起来朴素，但它实际上在做一次\textbf{语义噪声削减}：它不理解文本含义，但它先把会干扰后续分块和提示词的格式噪声尽量抹平。

\begin{casebox}[本次案例中的真实输入]
\begin{itemize}
  \item 文件：武汉大学品牌声誉深度分析报告.pdf
  \item 文件大小：2893754 字节，约 2.89MB
  \item 抽取后的总文本长度：20862 字符
\end{itemize}
\end{casebox}

\section{为什么文本清洗不能太激进}

清洗过轻，噪声残留；清洗过重，原文结构损失。比如：

\begin{itemize}
  \item 如果把所有换行都变成空格，原本的章节边界会被抹平；
  \item 如果做太强的句法重写，原文措辞会被篡改；
  \item 如果过滤“重复句”，可能会误删报告中用于强调的关键判断。
\end{itemize}

这说明第一步的目标不是“生成更优美的文本”，而是“生成一个尽量忠实、但更适合后续分块和抽取的工作文本”。

\begin{takeawaybox}[本章核心结论]
PDF 到文本的过程，是把“适合人阅读的视觉对象”转成“适合系统分块和抽取的语义序列”。这一步如果做不好，后面所有层都会被污染。
\end{takeawaybox}

\chapter{第二步：为什么必须分块}

\section{分块不是技术细节，而是系统的时间复杂度开关}

一旦文本长度上来，后续所有操作都不再适合在整篇文本上一次完成。分块的核心作用是把一个长上下文问题，变成多个局部处理问题。它既影响抽取质量，也影响速度和成本。

本项目的默认切块函数采用按字符长度切分，并尝试优先在句末符号、空行或自然边界处截断。其参数只有两个：

\[
\text{chunk\_size} = s,\qquad \text{overlap} = o
\]

其中 \(s\) 决定每块能容纳多少上下文，\(o\) 决定前后块共享多少边界信息。

\section{为什么 overlap 必须存在}

如果完全不重叠，那么位于块边界附近的实体关系可能被拦腰截断。一个典型例子是：主体在上一块出现，关系事实在下一块出现，模型就可能无法把两者连起来。

因此 overlap 的作用是：\textbf{用少量冗余换取边界连续性}。但 overlap 太大又会带来重复抽取、重复写入和更高的处理延迟。

\section{本次案例为什么只切成了 4 个 chunk}

本案例文本长度为 20862 字符，配置是：

\[
\text{chunk\_size}=6500,\qquad \text{overlap}=650
\]

结果是只切成了 4 个 chunk。这个结果有明显的好处，也有明显的问题。

\paragraph{好处：}
块数少，意味着：
\begin{itemize}
  \item 写入 Graphiti 的 episode 数少；
  \item 远端 LLM 抽取调用次数少；
  \item 整体建图速度更快。
\end{itemize}

\paragraph{问题：}
块太粗，对一篇高密度的中文分析报告来说，单块内部会混入太多角色、平台和叙事层次，导致：
\begin{itemize}
  \item 细粒度事实边界变模糊；
  \item 局部检索的匹配词命中更差；
  \item 报告阶段“按问题召回相关事实”的质量下降。
\end{itemize}

\begin{warningbox}[本次案例中已经出现的信号]
报告阶段本地搜索共调用 21 次，其中 18 次返回 0 条相关事实。这并不意味着图谱不存在，而更像是“图谱存在，但按当前问题措辞难以召回到足够相关的局部事实”。这正是大 chunk 方案的典型副作用。
\end{warningbox}

\begin{figure}[H]
\centering
\begin{tikzpicture}
\begin{axis}[
  width=0.88\textwidth,
  height=7cm,
  xlabel={chunk 粒度（由小到大）},
  ylabel={相对水平},
  xmin=0.8,xmax=3.2,
  ymin=0,ymax=1.05,
  xtick={1,2,3},
  xticklabels={较细（2500级）,中等（4000级）,较粗（6500级）},
  legend style={at={(0.02,0.98)},anchor=north west,draw=none,fill=none},
  grid=major
]
\addplot[smooth,thick,DeepBlue,mark=*] coordinates {(1,0.35) (2,0.65) (3,0.92)};
\addplot[smooth,thick,AccentGreen,mark=square*] coordinates {(1,0.88) (2,0.73) (3,0.42)};
\legend{速度,检索命中潜力}
\end{axis}
\end{tikzpicture}
\caption{chunk 越大通常越快，但对密集分析型中文报告，检索命中潜力往往下降}
\end{figure}

\begin{takeawaybox}[本章核心结论]
分块不是“为了适配模型上下文”这么简单，它本质上是系统在局部性、连续性和成本之间做的第一次全局权衡。
\end{takeawaybox}

\chapter{第三步：本体生成，本质上在做什么}

\section{本体不是标签集合，而是系统的观察框架}

本体生成的输入是两样东西：

\begin{enumerate}
  \item 文档文本；
  \item 你想问的问题，也就是模拟需求。
\end{enumerate}

这意味着本体不是从文本中“客观提取”出来的固定真相，而是\textbf{围绕任务构造的观察框架}。系统会问：如果我要模拟社交媒体舆情，那我到底需要哪些会发声的主体、哪些关系类型，以及它们之间怎样连线才有意义？

\section{为什么这里必须先有问题，再有本体}

如果问题是“总结武汉大学这份报告说了什么”，那么本体可能更偏向报告章节、指标、结论。

但我们的实际问题是：

\begin{quote}
如果武汉大学发布了对肖某某处分的撤销公告，大众舆情走向会怎么样？
\end{quote}

这使得本体必须偏向“谁会说话、谁会转发、谁会对谁产生影响”的社会互动结构，而不是偏向文档结构本身。

\section{本项目的本体生成约束}

本项目的本体生成提示词给了非常强的结构约束：

\begin{itemize}
  \item 必须正好输出 10 个实体类型；
  \item 最后两个必须是兜底类型：\code{Person} 和 \code{Organization}；
  \item 实体必须是真实世界中可以发声的主体，而不是抽象概念；
  \item 关系类型要服务于社交媒体舆论模拟。
\end{itemize}

这段约束极其重要，因为它回答了一个核心问题：\textbf{系统并不是什么都往图里塞，而是只允许“对后续仿真有用的主体”进入中间世界模型。}

\section{本次案例实际生成了什么}

本次案例中，本体生成阶段得到的 10 个实体类型包括：

\begin{itemize}
  \item \code{University}
  \item \code{GovernmentAgency}
  \item \code{MediaOutlet}
  \item \code{SocialMediaPlatform}
  \item \code{Student}
  \item \code{Expert}
  \item \code{AlumniGroup}
  \item \code{ParentCommunity}
  \item \code{Person}
  \item \code{Organization}
\end{itemize}

从第一性原理看，这个结果说明系统已经把原始 PDF 重新看成一个“舆情主体生态”，而不是一篇普通报告。

\begin{definitionbox}[本体层的真正作用]
本体层的真正作用不是描述世界，而是规定系统将以什么粒度、什么视角、什么可操作对象去重建世界。
\end{definitionbox}

\begin{takeawaybox}[本章核心结论]
本体生成是在回答“为了让这个问题可被模拟，我需要承认世界里有哪些会发声的对象，以及它们之间可能存在什么关系”。这是结构化建模的起点。
\end{takeawaybox}

\chapter{第四步：图谱构建，本质上在做什么}

\section{图谱是世界模型，不是数据库截图}

当 chunk 和本体都准备好之后，系统开始进入真正的“中间世界模型构造”阶段。

图谱构建过程可以概括为六步：

\begin{enumerate}
  \item 创建图谱容器；
  \item 设置本体；
  \item 对文本做 chunk 切分；
  \item 分批把 chunk 作为 episode 写入；
  \item 等待 Graphiti 处理；
  \item 读取节点和边统计结果。
\end{enumerate}

这六步不是机械流水线，而是在做一件非常具体的事：\textbf{把文本叙述压缩成节点、边和事实，交给图结构存储。}

\section{Graphiti Sidecar 在做什么}

本案例使用的是本地 Graphiti Sidecar。它有两个关键特点：

\begin{itemize}
  \item embedding 在本地，使用 \code{BAAI/bge-m3}；
  \item 信息抽取使用远端 OpenAI-compatible 生成模型，本次为 \code{gpt-5.4-mini}。
\end{itemize}

因此，图谱构建实际上是“本地向量化 + 本地图存储 + 远端结构抽取”的混合过程。它不是纯云，也不是纯本地。

\section{为什么图谱是比 chunk 更强的中间表示}

chunk 只是切开的文本片段，它仍然是线性的。

图谱则把文本重排为：

\begin{itemize}
  \item 节点：谁存在；
  \item 边：谁和谁有关；
  \item 边事实：这种关系在文本中被如何描述。
\end{itemize}

一旦进入图谱层，系统就可以做三件 chunk 难以高效完成的事：

\begin{enumerate}
  \item 按主体做聚合；
  \item 按关系做路径追踪；
  \item 按问题做局部检索。
\end{enumerate}

\section{本次案例图谱结果意味着什么}

本次运行最终得到：

\[
4\ \text{chunks},\quad 64\ \text{nodes},\quad 80\ \text{edges}
\]

这说明系统并没有把 20862 字符按句子级原样复制进来，而是压缩成了一个中等规模的关系网络。

\begin{casebox}[本次案例中的图谱意义]
64 个节点并不表示“只有 64 个事实”，而表示系统认为在这个舆情世界中，存在 64 个值得被识别和复用的主体或实体节点。80 条边则表示这些主体之间至少存在 80 条可用于解释和检索的显式关系。
\end{casebox}

\section{为什么建图速度会慢}

从代码流程看，真正耗时的不是单纯的 Kuzu 写盘，而是：

\begin{itemize}
  \item 把 chunk 发送给 Graphiti Sidecar；
  \item Sidecar 使用 LLM 做结构抽取；
  \item 抽取结果再转成图结构落到本地。
\end{itemize}

所以建图时间不是磁盘 I/O 的函数，而更接近：

\[
T_{\text{build}} \approx T_{\text{remote LLM extraction}} + T_{\text{embedding}} + T_{\text{graph write}}
\]

其中第一项通常是主导项。

\begin{takeawaybox}[本章核心结论]
图谱层的价值在于把线性文本变成可索引、可聚合、可局部访问的世界模型；建图层的成本则主要来自结构抽取，而不是图数据库本身。
\end{takeawaybox}

\chapter{第五步：为什么图谱还不够，必须继续生成 Agent}

\section{图谱里的实体不等于可发声主体}

图谱里有节点，并不意味着这些节点已经具备仿真能力。图谱的节点更像“世界中的对象”，而仿真需要的是“世界中的参与者”。

因此系统需要继续做一步转换：从图谱中的实体节点里筛出那些能够承担社交媒体角色的主体，然后为它们补齐行为属性。

\section{本项目如何从图谱生成 Agent 配置}

系统会先读取图谱中的全部节点与边，再过滤出符合预定义实体类型的节点。随后，系统通过
\code{SimulationConfigGenerator} 分阶段生成：

\begin{enumerate}
  \item 时间配置；
  \item 事件配置；
  \item Agent 配置；
  \item 平台配置。
\end{enumerate}

其中 Agent 配置不是一次性全生成，而是按批次生成。本项目的参数是：

\[
\text{AGENTS\_PER\_BATCH} = 15
\]

这说明系统已经意识到“主体人格化”本身也是一个长上下文问题，不能粗暴一次完成。

\section{为什么要把节点人格化}

因为图谱回答的是“谁和谁有什么关系”，而仿真回答的是“他们在什么平台、什么时间、以什么身份、用什么口吻、带着什么倾向发声”。

这两个问题不是一回事。

\begin{definitionbox}[从实体到 Agent 的本质]
实体是结构化世界里的对象；Agent 是结构化世界里的行为载体。前者解决“世界里有什么”，后者解决“这些东西会怎么动”。
\end{definitionbox}

\section{本次案例里的真实结果}

本次案例从图谱中生成了 59 个 Agent 配置。状态文件中的推理摘要还明确写出，这 59 个 Agent 并不是机械复制，而是经过以下几步形成：

\begin{itemize}
  \item 根据事件性质生成 4 天的时间仿真配置；
  \item 根据主题设计热点话题与事件节奏；
  \item 生成 59 个 Agent 配置；
  \item 为 8 个初始帖子分配发布者。
\end{itemize}

这说明系统已经把“文档中的知识”转换成了“平台上的主体生态”。

\begin{takeawaybox}[本章核心结论]
图谱是静态世界模型，Agent 集合是动态行为世界的入口。只有完成这一步，系统才能真正回答反事实舆情问题。
\end{takeawaybox}

\chapter{第六步：社会仿真如何把结构化记忆变成动态舆情}

\section{什么叫“在图谱上做仿真”}

这一步最容易被误解。系统并不是在图谱上“跑图算法”，而是在一个具有平台规则、时间推进和多主体互动的仿真环境里，让 Agent 依据自己的配置和记忆进行发言、评论、点赞、转发或沉默。

于是，原来静态的结构化记忆开始转化为动态轨迹：

\[
\mathcal{G} + \mathcal{A} + \tau \longrightarrow \mathcal{S}
\]

其中 \(\tau\) 表示时间推进机制。

\section{本次案例的仿真配置}

这次案例使用的是双平台并行仿真：

\begin{itemize}
  \item 平台：Twitter + Reddit；
  \item 总模拟时长设定：96 小时；
  \item 但因速度约束，实际执行轮数被截断为 12 轮；
  \item 共 59 个 Agent 参与；
  \item 最终动作总数为 216，其中 Twitter 108 个，Reddit 108 个。
\end{itemize}

\section{为什么仿真不是“一次多轮问答”}

因为平台会引入新的约束：

\begin{itemize}
  \item 不是所有人每轮都发言；
  \item 发言有时间分布；
  \item 不同平台的传播机制不同；
  \item 初始帖子会对后续互动路径产生偏置；
  \item 沉默也是一种行为。
\end{itemize}

因此仿真比普通多轮问答更接近“在规则约束下的多主体轨迹生成”。

\section{这次仿真为什么看起来并不快}

真实日志中有一个非常值得教学使用的细节：Twitter 在 16:44:05 就发出了 8 条初始帖子，而 Reddit 直到 16:47:17 才发出 8 条初始帖子。这说明双平台虽然逻辑上并行，但在真实执行中并没有完全并满。

这告诉我们一个重要事实：\textbf{系统标称并行，并不等于性能上就真正并行。}如果某一侧在等待远端模型、等待资源或等待上下文准备，整体体感仍会表现为“卡顿后跳跃”。

\begin{casebox}[本次仿真的真实耗时]
\begin{itemize}
  \item Twitter 模拟循环：226.2 秒
  \item Reddit 模拟循环：41.5 秒
  \item 仿真总循环耗时：241.1 秒
  \item 从 API 时间戳计算的仿真区间：246.5 秒
\end{itemize}
\end{casebox}

\section{为什么仿真进程结束后还在}

这也是一个真实案例中的常见误判点。仿真脚本默认不会在模拟循环完成后立即退出，而是进入“等待命令模式”，以便后续接受 interview 命令。也就是说，\textbf{进程还活着不代表跑挂了，可能只是为了支持报告阶段继续采访 Agent。}

\begin{warningbox}[教学提醒]
很多人看到 Python 进程不退出，会误以为系统卡住。这个案例恰好说明：应先看状态接口和脚本设计，再判断是否异常。这里进程保活是设计行为，而不是故障。
\end{warningbox}

\begin{takeawaybox}[本章核心结论]
仿真层把静态图谱变成动态舆情轨迹，它不是“多问几次模型”，而是在一个带平台、时间和主体差异的规则世界中生成行为序列。
\end{takeawaybox}

\chapter{第七步：报告为什么不是直接摘要，而是二阶推理}

\section{报告层真正依赖的不是原文，而是两个世界}

到报告阶段时，系统同时握有两类资源：

\begin{enumerate}
  \item 图谱世界：谁和谁有关、有哪些事实、有哪些节点和边；
  \item 仿真世界：各类主体在给定反事实情境下如何行动、如何表达、如何分化。
\end{enumerate}

因此报告不是在“再读一遍 PDF”，而是在这两个世界之上进行二阶组织。

\section{ReportAgent 的工作方式}

本项目中的 ReportAgent 采用 ReACT 模式，每个章节都可以调用工具。它主要依赖四类工具中的三类核心能力：

\begin{itemize}
  \item \code{panorama\_search}：从图谱中广角搜索全貌；
  \item \code{quick\_search}：做局部快速检索；
  \item \code{interview\_agents}：采访仿真中的 Agent；
  \item \code{insight\_forge}：对上下文做更深层次的问题推进。
\end{itemize}

这意味着它的写作过程不是“先想好再输出”，而是“边搜索、边采访、边写作、边修正”。

\begin{figure}[H]
\centering
\begin{tikzpicture}[
  tri/.style={rounded rectangle, draw=DeepBlue, thick, fill=SoftBlue, minimum width=3.3cm, minimum height=1.1cm, align=center},
  arrow/.style={-Latex, thick, DeepBlue}
]
\node[tri] (search) at (0,0) {图谱检索\\事实与结构};
\node[tri] (interview) at (4.9,0) {Agent 采访\\主观反应与立场};
\node[tri] (write) at (2.45,-3.1) {章节写作\\归纳、比较、判断};

\draw[arrow] (search) -- node[above]{补充事实} (interview);
\draw[arrow] (interview) -- node[right]{补充视角} (write);
\draw[arrow] (write) -- node[left]{提出新问题} (search);
\end{tikzpicture}
\caption{报告生成的本质是“检索 + 采访 + 写作”的闭环，而不是单次生成}
\end{figure}

\section{为什么报告阶段的搜索质量如此重要}

如果搜索能精准召回局部事实，报告就能建立在“结构化证据”上；如果搜索召回差，报告就会被迫更依赖 Agent 采访和模型归纳，从而：

\begin{itemize}
  \item 事实性变弱；
  \item 对问题措辞更敏感；
  \item 更容易出现“大方向对，但局部支撑薄”的现象。
\end{itemize}

本次案例中，报告阶段的本地搜索 21 次里有 18 次命中 0 条事实，这就是一个极强的教学信号：\textbf{速度调优不能只看前半段，必须看最终写作阶段的召回代价。}

\section{这次报告最后输出了什么}

本次报告最后得出的核心判断是：撤销公告不会直接平息舆情，而更可能把讨论从“处分争议”转成“撤销争议”，并进一步外溢到程序正义、公信力、高校治理和性别议题。

这个结论之所以有教学意义，不是因为它“猜得像”，而是因为它展示了系统的生成逻辑：

\begin{enumerate}
  \item 先从图谱层抽取已有舆情结构；
  \item 再从仿真主体中获取多角色反应；
  \item 最后把结构和视角拼成趋势叙事。
\end{enumerate}

\begin{takeawaybox}[本章核心结论]
报告层是二阶推理层。它消费的不是 PDF 原文本身，而是前面几层构造出的结构世界和行为世界。
\end{takeawaybox}

\chapter{武汉大学案例的真实运行复盘}

\section{时间分布}

本次案例的真实时间数据如下：

\begin{itemize}
  \item 文本抽取与本体阶段：15.5 秒
  \item 图谱构建阶段：170.2 秒
  \item 仿真准备阶段：185.2 秒
  \item 双平台仿真阶段：246.5 秒
  \item 报告生成阶段：301.9 秒
  \item 项目创建到报告完成全链路：939.3 秒
\end{itemize}

\begin{figure}[H]
\centering
\begin{tikzpicture}
\begin{axis}[
  width=0.9\textwidth,
  height=8cm,
  xbar,
  xmin=0,
  xmax=340,
  bar width=10pt,
  xlabel={耗时（秒）},
  symbolic y coords={报告生成,双平台仿真,仿真准备,图谱构建,文本抽取与本体},
  ytick=data,
  nodes near coords,
  nodes near coords align={horizontal},
  grid=major,
  enlargelimits=0.2
]
\addplot[fill=DeepBlue!70] coordinates {
  (301.9,报告生成)
  (246.5,双平台仿真)
  (185.2,仿真准备)
  (170.2,图谱构建)
  (15.5,文本抽取与本体)
};
\end{axis}
\end{tikzpicture}
\caption{武汉大学案例真实运行的各阶段耗时分布}
\end{figure}

\section{规模分布}

\begin{table}[H]
\centering
\begin{tabularx}{\textwidth}{p{4.5cm}X}
\toprule
指标 & 结果 \\
\midrule
原始文本长度 & 20862 字符 \\
Chunk 数量 & 4 \\
图谱节点数 & 64 \\
图谱边数 & 80 \\
可用 Agent 数 & 59 \\
仿真轮数 & 12（由 96 小时截断而来） \\
总动作数 & 216 \\
报告章节数 & 4 \\
报告搜索调用数 & 21 \\
其中 0 命中搜索 & 18 \\
\bottomrule
\end{tabularx}
\caption{本次案例的关键规模指标}
\end{table}

\section{为什么全链路比“官方演示”慢很多}

这是用户在真实使用中最容易提出的问题之一。对比这次案例与一些官方界面演示图，最核心的区别不是“代码逻辑有错”，而是执行环境不同：

\begin{itemize}
  \item 本次是本地 Graphiti Sidecar + 本地 Kuzu + 远端 LLM；
  \item 报告阶段还要依赖真实 Agent 采访；
  \item 双平台仿真和报告都具有明显的多次远端模型调用。
\end{itemize}

因此，体感上的“卡顿”主要来自生成式推理链路，而不是本机 CPU 算不动。

\begin{takeawaybox}[本章核心结论]
当你把真实时间分布拉平来看，会发现本系统最贵的不是“文档解析”，而是“生成配置、运行仿真、产出报告”这三步。它们都与远端推理有强耦合。
\end{takeawaybox}

\chapter{速度、质量与瓶颈：为什么会慢，为什么会失真}

\section{真正的主瓶颈：远端大模型链路}

本次案例中，最有解释力的结论是：\textbf{系统的主瓶颈不是本地图数据库，也不是本地 embedding，而是远端大模型调用链路。}

理由有三：

\begin{enumerate}
  \item 建图、仿真准备、仿真和报告这三段都大量依赖生成式调用；
  \item 日志表现为“停一段，再跳一段”，更像等待网络与接口响应，而不是持续高占用计算；
  \item 仿真进程存在多条指向本地代理端口的连接痕迹，说明远端访问链路本身就是系统时延的一部分。
\end{enumerate}

\section{第二瓶颈：主体数与采访数}

59 个 Agent 对于教学案例已经足够丰富，但对速度并不友好。原因在于：

\begin{itemize}
  \item 生成 59 个 Agent 配置本身就要多轮推理；
  \item 双平台仿真意味着行为决策空间翻倍；
  \item 报告阶段多次采访 Agent，会继续消耗接口时延。
\end{itemize}

从第一性原理看，Agent 数越多，世界越像真的，但时间成本也越像真的。

\section{第三瓶颈：大 chunk 带来的检索退化}

6500 / 650 的设置让前半段更快，但它削弱了后半段。对报告来说，检索命中差的代价是：

\begin{itemize}
  \item 需要反复遍历全图；
  \item 很多问题找不到直接支撑事实；
  \item 报告只能更多依赖采访和总结。
\end{itemize}

这就是一个非常典型的系统级 trade-off：

\begin{quote}
前端建图越快，不代表后端写报告越快；局部召回越差，最终分析反而可能更慢。
\end{quote}

\section{本次运行有没有错误}

如果严格按状态接口和日志判断，本次运行\textbf{没有致命错误}：

\begin{itemize}
  \item 仿真状态返回 \code{error: null}；
  \item 报告状态返回 \code{error: null}；
  \item 图谱构建任务成功完成；
  \item 最终报告文件已落盘。
\end{itemize}

但存在两个非致命 warning：

\begin{itemize}
  \item 某些 Bert 权重未初始化；
  \item 某个 attention 路径退回到了 manual implementation。
\end{itemize}

这两个 warning 更像依赖库和模型兼容性信号，而不是本次流程失败原因。

\begin{takeawaybox}[本章核心结论]
本次案例最值得学习的，不是“有没有报错”，而是“系统为什么在没有致命错误的前提下仍然可以表现得慢，而且部分阶段质量下降”。这正是复杂 AI 系统的常态。
\end{takeawaybox}

\chapter{如何做参数调优：不是越快越好，而是找到系统平衡点}

\section{调优要先决定优化目标}

同一套系统，至少可能有三种不同的优化目标：

\begin{enumerate}
  \item \textbf{最短总耗时}：尽快从 PDF 跑到报告；
  \item \textbf{最高报告质量}：让检索、采访和写作尽量有支撑；
  \item \textbf{最好成本效率}：让质量损失在可接受范围内，但每次运行更便宜。
\end{enumerate}

这三个目标通常不能同时最优。

\section{对本案例最有价值的调优旋钮}

\subsection*{1. chunk 参数}
对高密度中文分析 PDF，建议重点比较：

\begin{itemize}
  \item 6500 / 650：前半段快，但后半段召回压力大；
  \item 4000 / 400：更可能在速度与召回之间取得平衡；
  \item 2500 / 250：召回细，但建图和抽取调用会显著增加。
\end{itemize}

\subsection*{2. Agent 数量}
59 个 Agent 适合展示复杂舆情生态；如果目标是更快地做业务验证，可以先压到 25--40 个。

\subsection*{3. 报告阶段工具调用}
报告生成不是越多工具越好。若搜索质量低，增加工具调用次数只会把慢放大。

\subsection*{4. 模型选择}
本次采用 \code{gpt-5.4-mini} 是速度优先思路下的选择。如果换成更强模型，可以承受更大 chunk，但未必会得到更快的全链路时间。

\section{推荐的下一步实验}

如果目标是把本案例继续向“可用产品”推进，而不是仅做一次演示，那么最值得做的是同题 A/B：

\begin{enumerate}
  \item 保持 \code{gpt-5.4-mini} 不变；
  \item 把 chunk 从 6500 / 650 改成 4000 / 400；
  \item 比较四组指标：
    \begin{itemize}
      \item 建图耗时；
      \item 图谱规模；
      \item 报告阶段检索命中率；
      \item 报告完成耗时。
    \end{itemize}
\end{enumerate}

这个实验之所以重要，是因为它能回答一个真正系统级的问题：\textbf{前半段变快，是否真的让最终答案更快、更稳、更有证据支撑。}

\begin{takeawaybox}[本章核心结论]
调优不应该只盯住单一步骤，而应围绕“最终报告质量与总耗时”做闭环评估。系统优化的单位不是函数，而是完整链路。
\end{takeawaybox}

\chapter{结语：什么叫“可被推演的中间世界模型”}

回到最初的问题：为什么系统能从一份 PDF 走到一份分析报告？

答案不是“大模型很强”，也不是“图数据库很强”，而是这套系统在中间搭了一座桥：

\begin{center}
\Large
原始文档 $\rightarrow$ 结构化世界 $\rightarrow$ 可发声主体 $\rightarrow$ 动态轨迹 $\rightarrow$ 可解释结论
\end{center}

这座桥里最关键的不是哪一个单点技术，而是中间世界模型的存在。因为只有当世界被结构化、主体被人格化、轨迹被模拟化之后，系统才有条件把“过去文本”转化为“未来判断”。

\begin{definitionbox}[最终定义]
所谓“从 PDF 到分析”的本质，是构造一个足够压缩、足够结构化、足够可检索、足够可被推演的中间世界模型，并在其上完成多主体反事实分析。
\end{definitionbox}

这也是为什么本案例特别适合作为教学样本。它不仅展示了系统如何成功，还展示了系统在哪里慢、在哪里丢信息、在哪里需要调优。对于产品、研究和方案设计来说，这比一次漂亮但不可解释的生成结果更有价值。

\backmatter

\appendix

\chapter{附录：本次案例的真实配置与日志摘要}

\section{关键配置摘要}

\begin{table}[H]
\centering
\small
\begin{tabularx}{\textwidth}{p{4.2cm}X}
\toprule
字段 & 值 \\
\midrule
\code{GRAPH\_BACKEND} & \code{graphiti} \\
\code{GRAPHITI\_SERVICE\_URL} & \code{http://127.0.0.1:8011} \\
\code{LLM\_MODEL\_NAME} & \code{gpt-5.4-mini} \\
\code{GRAPHITI\_EMBEDDING\_MODEL} & \code{BAAI/bge-m3} \\
\code{GRAPHITI\_BUILD\_BATCH\_SIZE} & \code{2} \\
\code{DEFAULT\_CHUNK\_SIZE} & \code{6500} \\
\code{DEFAULT\_CHUNK\_OVERLAP} & \code{650} \\
\bottomrule
\end{tabularx}
\end{table}

\section{关键时间戳摘要}

\begin{table}[H]
\centering
\small
\begin{tabularx}{\textwidth}{p{4cm}p{3.6cm}X}
\toprule
阶段 & 时间点 & 说明 \\
\midrule
项目创建 & 16:37:37 & 项目开始 \\
图谱任务创建 & 16:37:52 & 开始正式建图 \\
图谱完成 & 16:40:43 & 图谱阶段结束 \\
仿真准备完成 & 16:43:52 & 59 个 Agent 配置生成完成 \\
仿真启动 & 16:43:53 & 双平台并行仿真开始 \\
仿真完成 & 16:47:59 & 216 个动作执行完毕 \\
报告开始 & 16:48:14 & 进入报告大纲与章节生成 \\
报告完成 & 16:53:16 & 最终报告落盘 \\
\bottomrule
\end{tabularx}
\end{table}

\section{检索质量摘要}

\begin{itemize}
  \item 报告阶段搜索调用总数：21
  \item 0 命中次数：18
  \item 非 0 命中次数：3
  \item 平均命中：1.714 条事实
  \item 最大命中：20 条事实
\end{itemize}

\section{典型日志片段}

\begin{tcolorbox}[colback=SoftGray,colframe=black!30,title=仿真日志摘录]
\small\ttfamily\raggedright
[16:43:57] OASIS 双平台并行模拟\par
[16:43:57] [通用LLM] model=gpt-5.4-mini\par
[16:44:05] [Twitter] 已发布 8 条初始帖子\par
[16:47:17] [Reddit] 已发布 8 条初始帖子\par
[16:47:58] 模拟循环完成! 总耗时: 241.1 秒
\end{tcolorbox}

\begin{tcolorbox}[colback=SoftGray,colframe=black!30,title=报告日志摘录]
\small\ttfamily\raggedright
[16:48:14] 图谱搜索: graph\_id=mirofish\_62ade01a1fa44a8e\par
[16:48:14] 使用本地搜索\par
[16:48:14] 本地搜索完成: 找到 0 条相关事实\par
[16:48:22] 大纲规划完成: 4 个章节\par
[16:53:16] 报告生成完成: report\_ffd323a8b87f
\end{tcolorbox}

\chapter{附录：把这份讲义复用到别的案例时，应该先改什么}

如果你想把这份讲义方法迁移到别的项目，优先改以下五项，而不是从头重写全部概念：

\begin{enumerate}
  \item 输入文档类型和长度；
  \item 任务问题，即反事实问题或分析问题；
  \item 本体约束，即允许哪些类型的主体进入世界模型；
  \item 平台与主体规模，即仿真生态的复杂度；
  \item 调优目标，是偏速度、偏质量还是偏成本。
\end{enumerate}

只要这五项被改清楚，整套“PDF $\rightarrow$ 世界模型 $\rightarrow$ 仿真 $\rightarrow$ 报告”的第一性原理仍然成立。

\end{document}
